\documentclass[a4paper,12pt]{article}

\usepackage[french]{babel}
\usepackage[T1]{fontenc} 
\usepackage[latin1]{inputenc}
%\usepackage{amsmath}

\author{R�mi Synave, Stefka Gueorguieva, Pascal Desbarats}
\title{Manuel de la biblioth�que POLYHEDRON}
\date{}

\begin{document}
\maketitle

\section{Utilisation}
Cette biblioth�que impl�mente les structures de donn�es n�cessaires au stockage des polyh�dres d'un maillage ainsi que des fonctions pour manipuler celles-ci.\\


\section{Structures de donn�es}
Les polyh�dres sont communs aux structure de donn�es \texttt{vf\_model} et \texttt{vef\_model}.\\
Les polyh�dres �tant les �l�ments de plus haut niveau, il n'y a pas d'adjacence � g�rer pour cette structure.\\ 
Seules les quatre faces sont stock�es.\\
\begin{verbatim}
typedef struct
{
  int fa1; \\premiere face
  int fa2; \\seconde face
  int fa3; \\troiseieme face
  int fa4; \\quatrieme face
}polyhedron;
\end{verbatim}

\section{Fonctions}
\textbullet \texttt{void polyhedron\_display(polyhedron p)}\\
Affichage des quatre faces d'un polyh�dre sous la forme :\\
\begin{verbatim}
POLYHEDRON : (fa1,fa2,fa3,fa4)
\end{verbatim}
~\\
\underline{Param�tres et type de retour :}\\
\texttt{p} : le polyh�dre.\\
\texttt{retour} : aucun.\\


\end{document}
